% ==============================================================================
% PENGATURAN FORMATTING & TATA LETAK SKRIPSI
% Sesuai Pedoman Penulisan Skripsi FST UNIBA Madura 2026
% ==============================================================================

% 1. SPASI & INDENTASI
\onehalfspacing                      % Spasi antar baris 1.5
\setlength{\parindent}{0.75cm}       % Indent alinea baru 0.75 cm (7 ketukan)
\setlength{\parskip}{0pt}            % Jarak antar paragraf 0pt (tanpa spasi kosong ekstra antar paragraf)
\raggedbottom                        % Cegah peregangan elastis spasi vertikal LaTeX (persis perilaku Word DOCX)

% Pengaturan Spasi Daftar (List: Enumerate & Itemize) tanpa spasi ekstra (identik format Word DOCX)
\setlist{
  nosep,                             % Hapus space before/after dan space ekstra antar butir
  topsep=0pt,
  partopsep=0pt,
  parsep=0pt,
  itemsep=0pt,
  leftmargin=0.75cm,                 % Indentasi sejajar alinea baru
  labelsep=0.25cm
}
\setlist[enumerate]{label=\arabic*., ref=\arabic*}
\setlist[itemize]{label=\textbullet}

% 2. PENAMAAN STRUKTUR DOKUMEN (BAHASA INDONESIA)
\renewcommand{\contentsname}{DAFTAR ISI}
\renewcommand{\listfigurename}{DAFTAR GAMBAR}
\renewcommand{\listtablename}{DAFTAR TABEL}
\renewcommand{\bibname}{DAFTAR PUSTAKA}
\renewcommand{\chaptername}{BAB}
\renewcommand{\appendixname}{LAMPIRAN}
\renewcommand{\figurename}{Gambar}
\renewcommand{\tablename}{Tabel}

% 3. FORMAT PENOMORAN PERSAMAAN, TABEL, DAN GAMBAR PER BAB
% Sesuai pedoman: Persamaan (3-18), Tabel 3.1, Gambar 3.1
\renewcommand{\thechapter}{\Roman{chapter}}
\renewcommand{\thesection}{\arabic{chapter}.\arabic{section}}
\renewcommand{\thesubsection}{\arabic{chapter}.\arabic{section}.\arabic{subsection}}
\renewcommand{\theequation}{\arabic{chapter}-\arabic{equation}}
\renewcommand{\thefigure}{\arabic{chapter}.\arabic{figure}}
\renewcommand{\thetable}{\arabic{chapter}.\arabic{table}}

% 4. FORMAT JUDUL BAB, SUB-BAB, DAN SUB-SUB-BAB (TITLESEC)
% Format Bab: BAB I <baris baru> JUDUL BAB (Kapital, Tebal, Rata Tengah; BAB X: 14pt, Judul: 12pt)
\titleformat{\chapter}[display]
  {\singlespacing\centering\bfseries}
  {\fontsize{14pt}{18pt}\selectfont\MakeUppercase{\chaptername\ \thechapter}}
  {0.5em}
  {\fontsize{12pt}{16pt}\selectfont\MakeUppercase}

% Spasi bab: tepat 3.0 cm dari sisi atas kertas (persis format Word proposal teman)
\titlespacing*{\chapter}{0pt}{-35pt}{1.5em}

% Format Sub-bab (Section): Tebal, Rata Kiri, Title Case (12pt)
\titleformat{\section}
  {\singlespacing\bfseries\fontsize{12pt}{16pt}\selectfont}
  {\thesection}
  {0.75em}
  {}
\titlespacing*{\section}{0pt}{1.0em}{0.2em}

% Format Sub-sub-bab (Subsection): Tebal, Rata Kiri (12pt)
\titleformat{\subsection}
  {\singlespacing\bfseries\fontsize{12pt}{16pt}\selectfont}
  {\thesubsection}
  {0.75em}
  {}
\titlespacing*{\subsection}{0pt}{0.8em}{0.2em}

% Otomatisasi eliminasi jeda spasi ekstra jika \subsection langsung mengikuti \section tanpa teks perantara
% (menghilangkan kesan adanya baris kosong atau space after paragraph antara 2.2 dengan 2.2.1 / 2.3 dengan 2.3.1)
\newif\ifdirectlyaftersection
\directlyaftersectionfalse

\pretocmd{\section}{\directlyaftersectiontrue}{}{}
\pretocmd{\subsection}{%
  \ifdirectlyaftersection
    \vspace{-0.85em}%
    \directlyaftersectionfalse
  \fi
}{}{}

\makeatletter
\everypar=\expandafter{\the\everypar \directlyaftersectionfalse}
\makeatother

% 5. LOGIKA PENOMORAN HALAMAN (FANCYHDR)
% Pedoman FST UNIBA 2026:
% - Bagian awal: Angka Romawi kecil di tengah bawah
% - Bagian utama: Angka Arab di kanan atas
% - Halaman pertama setiap bab baru: Angka Arab di tengah bawah

% Style halaman awal skripsi (Frontmatter)
\fancypagestyle{frontstyle}{
  \fancyhf{}
  \renewcommand{\headrulewidth}{0pt}
  \renewcommand{\footrulewidth}{0pt}
  \fancyfoot[C]{\thepage}
}

% Style halaman isi utama (Mainmatter)
\fancypagestyle{mainstyle}{
  \fancyhf{}
  \renewcommand{\headrulewidth}{0pt}
  \renewcommand{\footrulewidth}{0pt}
  \fancyhead[R]{\thepage}
}

% Style 'plain' otomatis dipanggil pada halaman pertama \chapter
\fancypagestyle{plain}{
  \fancyhf{}
  \renewcommand{\headrulewidth}{0pt}
  \renewcommand{\footrulewidth}{0pt}
  \fancyfoot[C]{\thepage}
}

% 6. PENGATURAN SPASI & CAPTION OBJEK GAMBAR & TABEL
% Pengendalian spasi sebelum dan sesudah float/in-text float (menghilangkan jeda ekstra seperti space after)
\setlength{\floatsep}{4pt plus 1pt minus 1pt}
\setlength{\textfloatsep}{6pt plus 1pt minus 1pt}
\setlength{\intextsep}{0pt}
\setlength{\dblfloatsep}{4pt plus 1pt minus 1pt}
\setlength{\dbltextfloatsep}{6pt plus 1pt minus 1pt}

% Tabel: Judul di atas tabel, jarak dengan tabel rapat, tanpa spasi ekstra di atas judul
\captionsetup[table]{
  labelfont={bf},
  textfont={normalfont},
  justification=centering,
  labelsep=space,
  position=top,
  aboveskip=0pt,
  belowskip=4pt
}

% Gambar: Judul di bawah gambar, huruf label tebal, tanpa spasi ekstra di bawah judul
\captionsetup[figure]{
  labelfont={bf},
  textfont={normalfont},
  justification=centering,
  labelsep=space,
  position=bottom,
  aboveskip=4pt,
  belowskip=0pt
}

% 7. FORMAT KUTIPAN DAFTAR PUSTAKA APA
\AtBeginDocument{
  \renewcommand{\bibname}{DAFTAR PUSTAKA}
  \renewcommand{\refname}{DAFTAR PUSTAKA}
}
\setlength{\bibhang}{0.75cm} % Hanging indent 0.75 cm
\setlength{\bibsep}{6.9pt}   % Jarak jeda antar-entri pustaka (identik dengan after=138 twips / 6.9 pt pada acuan Word)

% 8. PENYOROTAN KODE PROGRAM (LISTINGS)
\definecolor{codegray}{rgb}{0.5,0.5,0.5}
\definecolor{backcolour}{rgb}{0.97,0.97,0.97}
\definecolor{primaryblue}{rgb}{0.0,0.2,0.6}
\definecolor{stringgreen}{rgb}{0.1,0.5,0.2}

\lstdefinestyle{skripsicode}{
  backgroundcolor=\color{backcolour},
  commentstyle=\color{codegray}\itshape,
  keywordstyle=\color{primaryblue}\bfseries,
  numberstyle=\tiny\color{codegray},
  stringstyle=\color{stringgreen},
  basicstyle=\ttfamily\footnotesize\singlespacing,
  breakatwhitespace=false,
  breaklines=true,
  captionpos=b,
  keepspaces=true,
  numbers=left,
  numbersep=7pt,
  showspaces=false,
  showstringspaces=false,
  showtabs=false,
  tabsize=2,
  frame=single,
  rulecolor=\color{lightgray}
}
\lstset{style=skripsicode}

% 9. HYPERLINK SETUP (PDF METADATA & LINK COLORING)
\hypersetup{
  colorlinks=true,
  linkcolor=black,
  citecolor=black,
  urlcolor=blue,
  pdfborder={0 0 0}
}

% 10. FORMAT DAFTAR ISI, DAFTAR GAMBAR, DAN DAFTAR TABEL (TITLETOC)
% Format Titik-Titik Rapat Bersambung Penuh Hingga Nomor Halaman & Hanging Indent
\contentsmargin[1.5em]{1.2em}

% Format Bab / Bagian Utama (BAB I PENDAHULUAN ...................... 1)
% Untuk entri tanpa nomor (frontmatter, daftar pustaka, lampiran): tercetak biasa tebal dengan titik rapat
\titlecontents{chapter}[0pt]
  {\vspace{0.35em}\bfseries}
  {BAB \thecontentslabel\ }
  {}
  {\titlerule*[0.28pc]{.}\contentspage}

% Format Sub-bab / Section (1.1 Latar Belakang ..................... 1)
% Indent label 0.8cm, teks judul & baris lipat mulai pada 1.6cm (hanging indent)
\titlecontents{section}[1.6cm]
  {\normalfont}
  {\contentslabel[\thecontentslabel]{0.8cm}}
  {}
  {\titlerule*[0.28pc]{.}\contentspage}

% Format Sub-sub-bab / Subsection (2.2.1 Sistem .................... 6)
% Indent label 1.6cm, teks judul & baris lipat mulai pada 2.8cm (hanging indent)
\titlecontents{subsection}[2.8cm]
  {\normalfont}
  {\contentslabel[\thecontentslabel]{1.2cm}}
  {}
  {\titlerule*[0.28pc]{.}\contentspage}

% Format Gambar pada Daftar Gambar (Gambar 3.1 Topologi ........... 11)
% Label "Gambar X.X" lebar 2.5cm, baris lipat hanging indent sejajar pada 2.5cm
\titlecontents{figure}[2.5cm]
  {\normalfont}
  {\contentslabel[\textbf{Gambar \thecontentslabel}]{2.5cm}}
  {}
  {\titlerule*[0.28pc]{.}\contentspage}

% Format Tabel pada Daftar Tabel (Tabel 2.1 Matriks ................ 6)
% Label "Tabel X.X" lebar 2.2cm, baris lipat hanging indent sejajar pada 2.2cm
\titlecontents{table}[2.2cm]
  {\normalfont}
  {\contentslabel[\textbf{Tabel \thecontentslabel}]{2.2cm}}
  {}
  {\titlerule*[0.28pc]{.}\contentspage}


